\chapter{Lorentz 群与有限维场表示}
\label{chap:lorentz-fields}

\chapref{chap:rotation-group}讨论的 $SO(3)$ 只改变空间方向，时间仍是独立参数。狭义相对论把时间与空间合成四维时空，允许不同惯性系之间作洛伦兹推动；相应的连续时空对称性由正常正时 Lorentz 群 $SO^+(3,1)$ 描述。它不是紧群，因而场的有限维表示一般不可能同时保持正定内积；这正是把旋转表示论直接照搬过来时首先遇到的困难。

本章采用 $c=\hbar=1$ 与度规 $\eta=\operatorname{diag}(1,-1,-1,-1)$，先由间隔不变性推出群条件和 Lie 代数，再把复化代数分成两个彼此对易的 $\mathfrak{su}(2)$ 因子。$(j_L,j_R)$ 由此分类标量、Weyl 旋量、四矢量和 Dirac 旋量。单粒子质量与自旋不是这套有限维标签；它们留到下一章由 Poincar\'e 群的幺正表示给出。
\section{闵可夫斯基时空与洛伦兹群}

四矢量写作 $x^\mu=(t,\vb{x})$，闵可夫斯基内积为
\[
x\cdot y=\eta_{\mu\nu}x^\mu y^\nu=x^0y^0-\vb{x}\cdot\vb{y} .
\]
洛伦兹群是保持这一双线性型的实线性变换群，
\[
O(3,1)=\{\Lambda\in GL(4,\mathbb R)\mid
\Lambda^{\mathsf T}\eta\Lambda=\eta\}.
\]
由行列式和 $\Lambda^0{}_0$ 的符号，$O(3,1)$ 分为四个连通分支。包含单位元的分支满足 $\det\Lambda=1$、$\Lambda^0{}_0\geq1$，记为 $SO^+(3,1)$。空间反演 $P=\operatorname{diag}(1,-1,-1,-1)$ 和时间反演 $T=\operatorname{diag}(-1,1,1,1)$ 不在这个连通分支中，必须作为分立对称另行加入。

沿 $x$ 轴、快度为 $\zeta$ 的推进为
\[
\Lambda_x(\zeta)=
\begin{pmatrix}
\cosh\zeta&-\sinh\zeta&0&0\\
-\sinh\zeta&\cosh\zeta&0&0\\
0&0&1&0\\0&0&0&1
\end{pmatrix},
\qquad v=\tanh\zeta .
\]
转角是周期变量，而 $\zeta\in\mathbb R$ 无界；因此 $SO^+(3,1)$ 非紧。非紧性有一个直接后果：除平凡表示外，它没有非平凡有限维幺正表示。原因是有限维表示中推进生成元 $K_i$ 必须取反厄米矩阵（与 $J_i$ 满足 $[K_i,K_j]=-\ii\epsilon_{ijk}J_k$），无法与全部生成元同时厄米化。场仍可处在有限维但非幺正的 Lorentz 表示中；量子态所需的幺正性则由无限维的 Poincar\'e 表示实现。

\section{洛伦兹代数与有限维场表示}

令 $M_{\mu\nu}=-M_{\nu\mu}$ 为生成元，并定义
\[
J_i=\frac12\epsilon_{ijk}M_{jk},\qquad K_i=M_{0i}.
\]
$J_i$ 生成空间转动，$K_i$ 生成推进。取生成元为厄米算符的约定时，洛伦兹代数为
\begin{equation}
[J_i,J_j]=\ii\epsilon_{ijk}J_k,\qquad
[J_i,K_j]=\ii\epsilon_{ijk}K_k,\qquad
[K_i,K_j]=-\ii\epsilon_{ijk}J_k .
\label{eq:lorentz-algebra}
\end{equation}
最后一个负号把双曲推进与欧氏转动区分开来；两个不共线推进的交换子是一个转动生成元。

在复化代数中引入
\begin{equation}
A_i\equiv\frac12(J_i+\ii K_i),\qquad
B_i\equiv\frac12(J_i-\ii K_i).
\label{eq:complex-rotation-boost}
\end{equation}
代入\eqnrefs{eq:lorentz-algebra} 得
\begin{equation}
[A_i,A_j]=\ii\epsilon_{ijk}A_k,\qquad
[B_i,B_j]=\ii\epsilon_{ijk}B_k,\qquad
[A_i,B_j]=0,
\label{eq:so31-complex-decomposition}
\end{equation}
即 $\mathfrak{so}(3,1)_{\mathbb C}\cong\mathfrak{sl}(2,\mathbb C)\oplus\mathfrak{sl}(2,\mathbb C)$。反变换为
\[
J_i=A_i+B_i,\qquad K_i=-\ii(A_i-B_i).
\]
真实 Lorentz 代数不是两个独立的实 $\mathfrak{su}(2)$；复共轭交换两因子，
\begin{equation}
\overline{A_i}=B_i,\qquad \overline{B_i}=A_i,
\label{eq:real-form-conjugation}
\end{equation}
这才指定了 $\mathfrak{so}(3,1)$ 的实形式。

有限维不可约复表示因此由一对半整数 $(j_L,j_R)$ 标记，表示空间为 $V_{j_L}\otimes V_{j_R}$，维数为 $(2j_L+1)(2j_R+1)$。在该空间中
\[
\rho(J_i)=J_i^{(j_L)}\otimes I+I\otimes J_i^{(j_R)},\qquad
\rho(K_i)=-\ii\bigl(J_i^{(j_L)}\otimes I-I\otimes J_i^{(j_R)}\bigr).
\]
常用表示包括标量 $(0,0)$、左右 Weyl 旋量 $(\tfrac12,0)$ 与 $(0,\tfrac12)$、四矢量 $(\tfrac12,\tfrac12)$，以及自对偶和反自对偶二形式 $(1,0)$、$(0,1)$。

宇称保持转动而反转推进，时间反演在量子理论中是反幺正算符，作用为
\begin{align*}
P J_iP^{-1}&=J_i,& P K_iP^{-1}&=-K_i,
\\
T J_iT^{-1}&=-J_i,& T K_iT^{-1}&=K_i,
\qquad T\ii T^{-1}=-\ii.
\end{align*}
因此宇称交换左右表示，
\[
P:(j_L,j_R)\longleftrightarrow(j_R,j_L).
\]
若理论要求宇称在单个场多重态内实现，就必须把 $(j_L,j_R)$ 与 $(j_R,j_L)$ 配对；Dirac 旋量正是 $(\tfrac12,0)\oplus(0,\tfrac12)$。

群元素可写成
\[
D^{(j_L,j_R)}(\boldsymbol\theta,\boldsymbol\zeta)
=\exp\!\left[-\ii\boldsymbol\theta\cdot\rho(\vb{J})
-\ii\boldsymbol\zeta\cdot\rho(\vb{K})\right].
\]
限制到转动子群时，两个角动量按 Clebsch--Gordan 规则耦合，
\[
(j_L,j_R)\big|_{SO(3)}
\cong\bigoplus_{j=|j_L-j_R|}^{j_L+j_R}j .
\]
复共轭与张量积分别满足
\begin{align*}
\overline{(j_L,j_R)}&\cong(j_R,j_L),
\\
(j_L,j_R)\otimes(k_L,k_R)
&\cong
\bigoplus_{a=|j_L-k_L|}^{j_L+k_L}
\bigoplus_{b=|j_R-k_R|}^{j_R+k_R}(a,b).
\end{align*}
这些公式直接决定场的分量怎样组成 Lorentz 标量、矢量和张量。

\begin{derivation}[复化分解]
由\eqnrefs{eq:complex-rotation-boost}与\eqnrefs{eq:lorentz-algebra}，
\begin{align*}
4[A_i,A_j]
&=[J_i,J_j]+\ii[J_i,K_j]+\ii[K_i,J_j]-[K_i,K_j]\\
&=2\ii\epsilon_{ijk}(J_k+\ii K_k)
=4\ii\epsilon_{ijk}A_k.
\end{align*}
同理 $[B_i,B_j]=\ii\epsilon_{ijk}B_k$，而交叉项逐项抵消，得到\eqnrefs{eq:so31-complex-decomposition}。这一步只在复化后成立；真实形式还需附加\eqnrefs{eq:real-form-conjugation}。
\end{derivation}

\section{Clifford 代数与旋量}

质量壳关系 $p^2=m^2$ 对四动量是二次的。若要构造对 $p_\mu$ 一阶的相对论波动算符，需要矩阵 $\gamma^\mu$ 满足
\[
\{\gamma^\mu,\gamma^\nu\}=2\eta^{\mu\nu}I_4 .
\]
这是复 Clifford 代数 $\mathrm{Cl}_{1,3}(\mathbb C)$ 的定义关系。令
\[
S^{\mu\nu}=\frac{\ii}{4}[\gamma^\mu,\gamma^\nu],
\qquad
S(\Lambda)=\exp\!\left(-\frac{\ii}{2}\omega_{\mu\nu}S^{\mu\nu}\right),
\]
其中 $\omega_{\mu\nu}=-\omega_{\nu\mu}$ 是 Lorentz 变换的反对称无穷小参数（共 6 个实参数：3 个转动角与 3 个快度）。Dirac 旋量表示由
\[
S(\Lambda)\gamma^\mu S(\Lambda)^{-1}=\Lambda^\mu{}_\nu\gamma^\nu
\]
定义；两边左乘 $S^{-1}$、右乘 $S$，Clifford 关系给出
\begin{equation}
S(\Lambda)^{-1}\gamma^\mu S(\Lambda)
=\Lambda^\mu{}_{\nu}\gamma^\nu.
\label{eq:gamma-lorentz}
\end{equation}
式 \eqnrefs{eq:gamma-lorentz} 即上面相似变换的逆形式。所以 $\slashed p=\gamma^\mu p_\mu$ 按相似变换协变，Dirac 方程 $(\slashed p-m)\Psi=0$ 与 Lorentz 变换相容。

手征算符 $\gamma^5=\ii\gamma^0\gamma^1\gamma^2\gamma^3$ 满足 $(\gamma^5)^2=I$，投影 $P_{L,R}=(I\mp\gamma^5)/2$ 把 Dirac 旋量分成左右 Weyl 分量。电荷共轭矩阵 $C$ 满足 $C\gamma^\mu C^{-1}=-(\gamma^\mu)^{\mathsf T}$；Majorana 条件要求旋量等于其电荷共轭 $\psi=\psi^c\equiv C\bar\psi^{\mathsf T}$。复 Clifford 代数决定矩阵结构；是否能施加该 Majorana 实条件还取决于实形式和时空签名，不能仅由复维数判断。

$4\times4$ 复矩阵由 $I,\gamma^5,\gamma^\mu,\gamma^\mu\gamma^5,\sigma^{\mu\nu}$ 张成。因此 Dirac 双线性按 Lorentz 表示分解为
\[
\begin{array}{c|c}
\text{双线性}&\text{Lorentz 类型}\\ \hline
\bar\Psi\Psi&(0,0)\text{ 标量}\\
\bar\Psi\gamma^5\Psi&(0,0)\text{ 赝标量}\\
\bar\Psi\gamma^\mu\Psi&(\tfrac12,\tfrac12)\text{ 矢量}\\
\bar\Psi\gamma^\mu\gamma^5\Psi&(\tfrac12,\tfrac12)\text{ 赝矢量}\\
\bar\Psi\sigma^{\mu\nu}\Psi&(1,0)\oplus(0,1)
\end{array}
\]
表中 $1+1+4+4+6=16$，与矩阵空间维数一致。这一完备性是 Fierz 重排的代数来源。
